\begin{center}
  \begin{tabular}{|c|c|l|} \hline
    \textbf{Subtask} & \textbf{Điểm} & \textbf{Giới hạn} \\ \hline
    1 & $10$ & $N=4$ \\ \hline
    2 & $10$ & $N\le 5$ \\ \hline
    3 & $10$ & $N\le 6$ \\ \hline
    4 & $10$ & $N\le 7$ \\ \hline
    5 & $10$ & $N\le 8$ \\ \hline
    6 & $10$ & $N\le 9$ \\ \hline
    7 & $10$ & $N\le 10$ \\ \hline
    8 & $30$ & Không có ràng buộc gì thêm \\ \hline
  \end{tabular}
\end{center}

\textbf{Cách tính điểm}

Giả sử test có tối đa $1$ điểm. Nếu trong bất kỳ lần gọi hàm \texttt{solve} nào, xâu $U$ bạn đưa ra không thỏa mãn với điều kiện của đề bài thì bạn sẽ được $0$ điểm cho toàn bộ test đó. Ngược lại, điểm của bạn sẽ được tính như sau:
\begin{itemize}
   \item Mỗi lần gọi hàm \texttt{solve} sẽ có điểm của lần gọi hàm đó (gọi là $score$). Xét lần gọi hàm tương ứng, gọi giá trị $P=Q+|U|$ là độ tối ưu kết quả của thí sinh, giá trị $J=Q+|U|$ là độ tối ưu kết quả của ban tổ chức ($Q$ là tổng số lần thực hiện truy vấn \texttt{ask}, $U$ là xâu đáp án trả về đã thỏa mãn điều kiện). Khi đó:
   \begin{itemize}
      \item Nếu $Q>10$ thì $score=0$.
      \item Nếu $N<P-J$ thì $score=0.4\times 0.9^{P-J-N-1}$.
      \item Nếu $4\le P-J\le N$ thì $score=0.4+0.4\times 0.8^{P-J-4}$.
      \item Nếu $0<P-J<4$ thì $score=0.8+0.2\times 0.8^{P-J}$.
      \item Nếu $P\le J$ thì $score=1$.
   \end{itemize}
   \item Điểm của một test sẽ bằng điểm (giá trị $score$) nhỏ nhất trong tất cả lần gọi hàm \texttt{solve}.
\end{itemize}

Điểm của mỗi subtask là điểm thấp nhất của một test trong subtask đó.